卷积神经网络的历史脉络
我们知道，所谓动物的“高级”特性，体现出来的是行为表象。更深层地，它们也会体现在大脑皮层的进化上。

在本节中，我们主要回顾了卷积神经网络的发展史，从历史的脉络中，我们可得出一个意象上的结论：深而洞察的专注可能要胜过广而肤浅的观察。这也从方法论上部分解释了“深度学习”的成功所在。

\subsection{肤浅的全面or局部的深入?CNN就是来回答这个问题的!}
BP神经网络的问题在哪里？

在前面的文章中，我们介绍了反向传播(Back Propagation，简称BP) 算法，在本质上，作为一种高效的优化算法（即帮助快速找到解的算法），BP算法加持的通常一种全连接神经网络。BP算法也有很多成功的应用，但只能适用于“浅层”网络（通常小于7层），因为“肤浅”，所以也就限制了它的特征表征能力，进而也就局限了它的应用范围。

为什么它难以“深刻”呢？

在很大程度上问题就出在它的“全连接”上。

难道“全连接”不好吗？

它更全面啊，难道全面反而是缺陷？

我们暂时不讨论这个问题，等读者朋友阅读完本专题的系列文章之后，答案自然就会了然于胸。

在本章，我们讨论一种应用范围更为广泛的网络---卷积神经网络(Convolutional Neural Network，简称CNN)，它在图像、语音识别等众多任务（比如GoogleNet、微软的ResNet等）上发挥神勇，近几年深度学习大放异彩，CNN可谓是功不可没，拔得头筹。

可为什么CNN能这么生猛呢？

答案还得从历史中追寻。著名人类学家费孝通先生曾指出[1]，
{\em 我们所谓的“当前”，其实包含着从“过去”历史中拔萃出来的投影和时间选择的积累。历史对于我们来说，并不是什么可有可无的点缀之饰物，而是实用的、不可或缺的前行之基础。}

下面我们就先聊聊卷积神经网络的历史脉络，希望能从中找到一点启发。在回顾历史之前，我们先尝试思考这样一个“看似题外话而实则不然”的问题：为什么几乎所有低级动物的双眼都是长在头部两侧？

12.2 眼在何方？路在何处？
的确，如果你仔细观察，低级动物的双眼大多都长在两侧。从进化论的角度来看，“物竞天择，适者生存”。

大自然既然这么安排，自然有它的道理。其中一个解释就是，那些低级动物正是因为这样的“造物安排”，所以能够同时看到上下左右前后等各个方向，从而就不存在视觉盲区。这确实是一种极为安全的配置，有了安全性，它们才能更好地在地球上生存。

可这样的配置又有什么局限呢？相比于低级动物，人的双眼可都是长在面部前方的。这样的配置，肯定不能全方位地观察周遭的一切，这岂不是很糟糕？但事实上是，只有人类进化成为这个地球上最为“高级”的动物。

有人是这样解释的（这个解释的意义可能更多的是来自于意象上的，而非生物学的，所以读者也无需死磕）：低级动物无死角的眼睛配置，虽然能够更全面的关注周围，但副作用却在于，它们没办法把自己的目光集中在某一处，自然也没有办法仔细、长期地观察某个点，于是它们也就不可能进化出深入思考的能力。而人类却因为眼睛的缺陷（接受了视野中的盲区）而能注视前方，从而能给出观察事物的深刻洞察（Insight）。“高级”动物就是这么“练就”出来了。

换句话说，肤浅的全面观察，有时候还不如局部的深入洞察。

可这，和我们今天的主题---卷积神经网络又有什么联系呢？当然有联系，这个联系体现在方法论层面，且听我慢慢道来。